\documentclass[12pt,a4paper]{article}
\usepackage{ucebnice}
\usepackage{qrcode}


\title{Logaritmické rovnice, nerovnice}
\author{Ondřej Škubala}
\date{\today}

\begin{document}

\maketitle
\newpage

\section{Úvod}

\subsection*{Součin uvnitř logaritmu}

Než začneme zapisovat a dokazovat obecné vzorce pro počítání s logaritmy, je vhodné
si připomenout, jak fungují mocniny se stejným základem. Vezmeme jednoduchý základ
$a=2$ a sestavíme tabulku hodnot pro několik vybraných exponentů:

\[
y = 2^x
\]

\begin{center}
\begin{tabular}{c|ccccccccc}
$x$     & 0 & 1 & 2 & 3 & 4  & 5  & 6  & 8   & 10   \\
\hline
$2^x$   & 1 & 2 & 4 & 8 & 16 & 32 & 64 & 256 & 1024 \\
\end{tabular}
\end{center}

Z tabulky snadno vyčteme, že
\[
16 = 2^4,\qquad 64 = 2^6,\qquad 1024 = 2^{10}.
\]

Zapisujeme tedy logaritmy:
\[
\log_2 16 = 4,\qquad \log_2 64 = 6,\qquad \log_2 1024 = 10.
\]

Protože však současně platí $16 \cdot 64 = 1024$, můžeme spojit všechny tři zápisy do jediné
rovnosti:
\[
\log_2 16 + \log_2 64 = 4 + 6 = 10 = \log_2 1024.
\]

Tato jednoduchá úvaha ukazuje, že logaritmus součinu lze přepsat jako součet logaritmů.

\begin{examplebox}
Z tabulky mocnin dvojky vyčteme
\[
16 = 2^4,\qquad 64 = 2^6,\qquad 1024 = 2^{10}.
\]
Proto platí
\[
\log_2 16 = 4,\qquad
\log_2 64 = 6,\qquad
\log_2 1024 = 10.
\]
A protože $16 \cdot 64 = 1024$, dostáváme
\[
\log_2 16 + \log_2 64 = 4 + 6 = 10 = \log_2 1024.
\]
Tedy konkrétní ukázku vzorce
\[
\log_2(16 \cdot 64) = \log_2 16 + \log_2 64.
\]
\end{examplebox}


Nyní si stejnou myšlenku ukážeme obecněji, bez konkrétních čísel. Představme si,
že máme dvě kladná čísla $r$ a $s$ a jejich logaritmy o základu $a$ (kde $a>0$, $a \neq 1$)
jsou $\log_a r$ a $\log_a s$. Podle definice logaritmu platí

\[
a^{\log_a r} = r
\quad\text{a}\quad
a^{\log_a s} = s.
\]

Pokud nyní obě rovnosti vynásobíme, dostaneme

\[
r \cdot s = a^{\log_a r} \cdot a^{\log_a s}.
\]

Součin dvou mocnin se stejným základem můžeme přepsat jako jednu mocninu, ve
které se exponenty sečtou:

\[
a^{\log_a r} \cdot a^{\log_a s} = a^{\log_a r + \log_a s}.
\]

Máme tedy
\[
r \cdot s = a^{\log_a r + \log_a s}.
\]

Na druhou stranu ale víme, že $r \cdot s$ je číslo, jehož logaritmus o základu $a$
je $\log_a (r \cdot s)$, tedy platí

\[
r \cdot s = a^{\log_a (r \cdot s)}.
\]

Tento zápis ukazuje, že číslo $r \cdot s$ lze vyjádřit jako mocninu se základem $a$
dvěma různými způsoby:

\[
a^{\log_a (r \cdot s)} = a^{\log_a r + \log_a s}.
\]

Protože je funkce $a^x$ prostá (každé hodnotě $x$ odpovídá právě jedna hodnota
$a^x$), musí být stejné také exponenty:

\[
\log_a (r \cdot s) = \log_a r + \log_a s.
\]

Tím jsme získali první základní vzorec pro práci s logaritmy.

\begin{definitionbox}
\textbf{Vzorec pro součin uvnitř logaritmu.}

Nechť $a>0$, $a\neq 1$ a nechť $r_1, r_2, \dots, r_n$ jsou kladná reálná čísla. Potom platí
\[
\log_a(r_1 r_2 r_3 \cdots r_n)
=
\log_a r_1 + \log_a r_2 + \log_a r_3 + \dots + \log_a r_n.
\]

Zvláštním případem je
\[
\log_a(r \cdot s) = \log_a r + \log_a s.
\]
\end{definitionbox}



\subsection*{Podíl uvnitř logaritmu}

Stejným způsobem, jakým jsme v předchozí části získali vztah pro logaritmus součinu,
můžeme odvodit i vzorec pro logaritmus podílu. Opět využijeme jednoduchý příklad
z tabulky mocnin dvojky.

Více již víme, že
\[
1024 = 2^{10}, \qquad 64 = 2^{6}, \qquad 16 = 2^{4}.
\]

Zároveň platí
\[
\frac{1024}{64} = 16.
\]

Zapisemeli logaritmy, dostáváme
\[
\log_2 1024 = 10,\qquad \log_2 64 = 6,\qquad \log_2 16 = 4,
\]
a okamžitě vidíme, že
\[
10 - 6 = 4.
\]

Tedy
\[
\log_2 1024 - \log_2 64 = \log_2 16.
\]

Protože $16 = \frac{1024}{64}$, můžeme výsledek zapsat i kompaktněji:
\[
\log_2 1024 - \log_2 64 = \log_2 \left( \frac{1024}{64} \right).
\]

Tato ukázka nám napovídá analogický vzorec ke vzorci pro součin: logaritmus podílu
je rozdílem logaritmů.


\begin{definitionbox}
\textbf{Vzorec pro podíl uvnitř logaritmu.}

Pro všechna $a>0$, $a\neq 1$ a pro každá dvě kladná reálná čísla $r, s$ platí
\[
\log_a\left(\frac{r}{s}\right)
=
\log_a r - \log_a s.
\]

Obecněji pro podíl více členů píšeme
\[
\log_a\!\left(\frac{r_1}{r_2 r_3 \cdots r_n}\right)
=
\log_a r_1 - \log_a r_2 - \log_a r_3 - \dots - \log_a r_n.
\]
\end{definitionbox}


\begin{notebox}
Existuje i podílový vzorec pro více než dva členy současně, ale jeho zápis je
nepraktický a v praxi zbytečný. Vystačíme si s opakovanou aplikací základního vztahu
\[
\log_a\left(\frac{r}{s}\right) = \log_a r - \log_a s.
\]
\end{notebox}



\subsection*{Logaritmus mocniny}

Po součinu a podílu přichází třetí základní úprava, která se týká výrazů, jež mají
v argumentu logaritmu mocninu. Je přirozené se ptát, zda lze logaritmus mocniny
převést na vhodný součet — a ukáže se, že ano, a to velmi elegantně.

Začněme jednoduchým příkladem.

\begin{examplebox}
Uvažujme výraz
\[
\log_2(2^3).
\]

Protože $2^3 = 2 \cdot 2 \cdot 2$, můžeme psát
\[
\log_2(2^3)
=
\log_2(2\cdot 2 \cdot 2).
\]

Použitím vzorce pro logaritmus součinu dostaneme
\[
\log_2(2\cdot 2 \cdot 2)
=
\log_2 2 + \log_2 2 + \log_2 2.
\]

Každý z těchto logaritmů má hodnotu $1$, takže
\[
\log_2(2^3) = 1 + 1 + 1 = 3 \cdot \log_2 2.
\]

Exponent $3$ se tedy přenesl před logaritmus jako násobek.
\end{examplebox}

Tento postup nestojí na žádné zvláštnosti konkrétního čísla $2$ ani exponentu $3$,
ale funguje obecně. Pokud číslo $r$ umocníme na $s$-tou, znamená to, že jej
vynásobíme samo sebou $s$-krát:
\[
r^s = \underbrace{r \cdot r \cdot \dots \cdot r}_{s\ \text{členů}}.
\]

A protože logaritmus převádí součin na součet, dostáváme opakovaným použitím
součinového pravidla
\[
\log_a(r^s)
=
\log_a r + \log_a r + \dots + \log_a r
=
s \cdot \log_a r.
\]

Tím získáváme třetí základní logaritmický vzorec.


\begin{definitionbox}
\textbf{Vzorec pro logaritmus mocniny.}

Pro všechna $a>0$, $a\neq 1$, libovolné kladné číslo $r$
a libovolné reálné $s$ platí:
\[
\log_a(r^s) = s \cdot \log_a(r).
\]

Tento vzorec říká, že exponent v argumentu logaritmu lze převést jako
násobek před logaritmus.
\end{definitionbox}



\subsection*{K čemu jsou logaritmické vzorce dobré?}

Dosud jsme se u exponenciálních rovnic setkávali převážně s „hezky vyrobenými“
příklady, kde šlo obě strany rovnice upravit na stejný základ a následně jen porovnat
exponenty, což vedlo k rychlému a čistému řešení. V praxi však často narazíme na
situace, kdy se základy na obou stranách rovnice nechtějí „sjednotit“ žádným rozumným
způsobem, a právě tam přichází ke slovu logaritmus a naše tři vzorce pro součin, podíl
a mocninu uvnitř logaritmu.

Podívejme se na rovnici
\[
3^{x+1} = 2^{3-2x},
\]
která na první pohled nevypadá nijak hrozivě, ale při pokusu o klasickou úpravu velmi
rychle narazíme na limity čistě mocninných postupů.

\begin{examplebox}
\textbf{Pokus o řešení bez logaritmu}

Začneme tím, že se pokusíme rovnici $3^{x+1} = 2^{3-2x}$ upravit pouze pomocí
mocninových pravidel. Přepíšeme obě strany tak, aby se exponenty lépe oddělily:

\[
3^{x+1} = 3^x \cdot 3,\qquad
2^{3-2x} = 2^3 \cdot 2^{-2x} = 8 \cdot 2^{-2x}.
\]

Rovnice tedy přejde na tvar
\[
3^x \cdot 3 = 8 \cdot 2^{-2x}.
\]

Obě strany vydělíme osmi a zároveň vynásobíme $2^{2x}$, abychom dostali všechny
proměnné do exponentu na jedné straně:
\[
\frac{3^x \cdot 3}{8} \cdot 2^{2x} = 1.
\]

Výrazy s $x$ můžeme spojit dohromady:
\[
\frac{3}{8} \cdot 3^x \cdot 2^{2x}
=
\frac{3}{8} \cdot (3 \cdot 4)^x
=
\frac{3}{8} \cdot 12^x.
\]

Rovnice se tedy zjednoduší na
\[
\frac{3}{8} \cdot 12^x = 1,
\]
odkud po vynásobení osmi třetinami dostaneme
\[
12^x = \frac{8}{3}.
\]

Tento tvar je sice přehledný a čistý, ale už nevede k jednoduchému „uhodnutí“ řešení:
neexistuje žádné pěkné přirozené číslo $x$, pro které by $12^x = \dfrac{8}{3}$ platilo.
V tuto chvíli nám samotná mocninová pravidla nestačí a potřebujeme silnější nástroj.
\end{examplebox}

Tímto silnějším nástrojem je právě logaritmus, který dokáže „stáhnout“ exponent dolů
a převést exponenciální rovnici na obyčejnou lineární rovnici, kterou již umíme snadno
vyřešit.

\begin{examplebox}
\textbf{Řešení rovnice pomocí logaritmu}

Vycházíme z upraveného tvaru
\[
12^x = \frac{8}{3},
\]
který jsme získali v předchozím příkladu. Na obě strany rovnice nyní vezmeme
logaritmus se stejným základem (na volbě základu nezáleží, může to být například
desítkový logaritmus, přirozený logaritmus nebo logaritmus o základu $2$):

\[
\log\left(12^x\right) = \log\left(\frac{8}{3}\right).
\]

Pomocí vzorce pro logaritmus mocniny vytáhneme exponent $x$ před logaritmus,
na pravé straně využijeme vzorec pro logaritmus podílu:
\[
x \cdot \log 12
=
\log 8 - \log 3.
\]

Nyní už jde jen o obyčejnou lineární rovnici v neznámé $x$, takže stačí vydělit
koeficientem u $x$:
\[
x = \frac{\log 8 - \log 3}{\log 12}.
\]


Vidíme tedy, že kombinací tří základních vzorců
\[
\log_a(rs) = \log_a r + \log_a s,\quad
\log_a\left(\frac{r}{s}\right) = \log_a r - \log_a s,\quad
\log_a(r^s) = s \log_a r
\]
umíme řešit i takové rovnice, kde čistě mocninné úpravy selhávají.
\end{examplebox}



\section{Rovnice s logaritmy}

Dosud jsme pracovali především s exponenciálními rovnicemi, v nichž byla neznámá
vystupující v exponentu, zatímco samotná rovnice byla zapsána pomocí mocnin. Nyní
se setkáme s rovnicemi, ve kterých se neznámá objevuje přímo uvnitř logaritmu, případně
dokonce v jeho základu, tedy s rovnicemi typu
\[
\log_2 x = 3,\qquad \log_x 8 = 3,
\]
a mnoha podobnými.

V rovnici $\log_2 x = 3$ hledáme takové $x>0$, pro které platí, že logaritmus o základu
$2$ má hodnotu $3$. Zapsáno pomocí exponenciálního tvaru to znamená, že hledáme
číslo $x$, pro které platí
\[
2^3 = x,
\]
tedy řešením je $x = 8$. Podobně u rovnice $\log_x 8 = 3$ hledáme takový základ $x>0$,
$x \neq 1$, aby logaritmus čísla $8$ o tomto základu byl roven $3$, což je ekvivalentní
zadání
\[
x^3 = 8,
\]
odkud vyplývá $x = 2$ (při splnění podmínek pro základ logaritmu).

\subsection*{Grafická představa řešení rovnice s logaritmem}

Rovnice $\log_2 x = 3$ má velmi názornou grafickou interpretaci. Uvažujme graf
funkce
\[
y = \log_2 x,
\]
tedy logaritmické funkce o základu $2$. Řešením rovnice $\log_2 x = 3$ jsou takové body
na grafu, které mají souřadnici $y = 3$. Geometricky tedy hledáme průsečík grafu
funkce $y = \log_2 x$ s vodorovnou přímkou $y = 3$.

\begin{center}
\begin{tikzpicture}[scale=1]
\begin{axis}[
    axis lines=middle,
    xmin=-1, xmax=10,
    ymin=-2, ymax=4.5,
    xlabel={$x$}, ylabel={$y$},
    samples=200,
    domain=0.1:10,
    smooth,
    ticks=none,
    width=11cm, height=7cm,
    title={\small Graf funkce $y = \log_2 x$}
]
% log_2 x = ln(x)/ln(2)
\addplot[thick] {ln(x)/ln(2)};

% vodorovná přímka y=3
\addplot[dashed] coordinates {(0.1,3) (10,3)};

% vertikální přímka x=8
\addplot[dashed] coordinates {(8,-2) (8,4.5)};

% Bod (8,0)
\addplot[only marks, mark=|] coordinates {(8,0)};
\node[below right] at (axis cs:8,0) {\footnotesize $8$};
% Bod (0,3)
\addplot[only marks, mark=|] coordinates {(0,3)};
\node[left] at (axis cs:0,3) {\footnotesize $3$};

% průsečík (8,3)
\addplot[only marks, mark=*] coordinates {(8,3)};
\node[above right] at (axis cs:8,3) {$[8;3]$};
\end{axis}
\end{tikzpicture}
\end{center}

Na obrázku je dobře vidět, že graf funkce $y = \log_2 x$ protíná přímku $y = 3$ v bodě
$[8;3]$, což odpovídá řešení rovnice $\log_2 x = 3$, tedy kořeni $x = 8$. Souřadnice
x tohoto průsečíku nám udává řešení rovnice, zatímco souřadnice y představuje hodnotu
logaritmu, kterou rovnice předepisuje.

\subsection*{Zlogaritmování a odlogaritmování rovnice}

Obdobně jako u exponenciálních rovnic jsme se snažili dostat obě strany do tvaru
\[
a^{\text{výraz}_1} = a^{\text{výraz}_2}
\quad\Longleftrightarrow\quad
\text{výraz}_1 = \text{výraz}_2,
\]
budeme nyní u rovnic s logaritmy usilovat o úpravu do tvaru
\[
\log_a(\text{výraz}_1) = \log_a(\text{výraz}_2)
\quad\Longleftrightarrow\quad
\text{výraz}_1 = \text{výraz}_2,
\]
kde $a>0$, $a\neq 1$ a oba výrazy musí mít kladné hodnoty, aby byl logaritmus definován.

Tato úprava funguje proto, že logaritmická funkce o pevném základu $a$ je prostá,
tedy každé hodnotě argumentu odpovídá právě jedna hodnota logaritmu a naopak.
Pokud jsou tedy logaritmy dvou kladných čísel stejné, musí být stejná i tato čísla.

V praxi budeme často postupovat dvěma opačnými směry:

buď rovnost vhodně \textbf{zlogaritmujeme}, tedy z rovnice typu
  \[
  a^{\text{výraz}} = \text{číslo}
  \]
  přejdeme na rovnost logaritmů, kterou pak upravíme pomocí logaritmických vzorců;

nebo naopak \textbf{odlogaritmujeme}, tedy z rovnice ve tvaru
  \[
  \log_a(\text{výraz}) = \text{číslo}
  \]
  přejdeme na ekvivalentní exponenciální rovnici, kterou vyřešíme pomocí mocnin.

\begin{notebox}
V této kapitole budeme používat neformální pojmy \textbf{zlogaritmování}
(„vezmeme logaritmus z obou stran rovnice“) a \textbf{odlogaritmování} („převedeme
rovnici s logaritmem zpět na exponenciální tvar“). V obou případech jde o využití
základní ekvivalence
\[
\log_a b = c \quad\Longleftrightarrow\quad a^c = b.
\]
\end{notebox}

\begin{examplebox}
\textbf{Řešte rovnici:}
\[
\log_2(x-2) = 4.
\]
Nejprve určíme podmínky pro výraz uvnitř logaritmu. Zápis $\log_2(x-2)$ je definován
právě tehdy, když argument logaritmu je kladný:
\[
x - 2 > 0 \quad\Longrightarrow\quad x > 2.
\]

Nyní rovnici odlogaritmujeme, tedy převedeme ji do exponenciálního tvaru pomocí
ekvivalence $\log_2 y = 4 \Longleftrightarrow y = 2^4$:
\begin{align*}
\log_2(x-2) &= 4 \\
x - 2 &= 2^4 \\
x - 2 &= 16 \\
x &= 18.
\end{align*}

Zkontrolujeme, zda nalezené řešení splňuje podmínku $x>2$, což je zřejmé. Proto je
\[
\text{Řešení: } K= \{18\}.
\]
\end{examplebox}

\begin{examplebox}
\textbf{Řešte rovnici:}
\[
3 \cdot \log(3x+1) = 6.
\]


Nejprve určujeme podmínky. Výraz $\log(3x+1)$ je definován tehdy, když argument
logaritmu je kladný:
\[
3x + 1 > 0 \quad\Longrightarrow\quad x > -\frac{1}{3}.
\]

Nyní rovnici upravíme. Nejprve obě strany vydělíme třemi:
\[
3 \cdot \log(3x+1) = 6
\quad\Longrightarrow\quad
\log(3x+1) = 2.
\]

Aplikujeme odlogaritmování (tj. ekvivalenci $\log y = 2 \Longleftrightarrow y = 10^2$,
protože zde používáme dekadický logaritmus):
\[
\log(3x+1) = 2
\quad\Longleftrightarrow\quad
3x + 1 = 10^{2}.
\]

Pokračujeme jako u lineární rovnice:
\begin{align*}
3x + 1 &= 100 \\
3x &= 99 \\
x &= 33.
\end{align*}

Zkontrolujeme splnění podmínky $x > -\tfrac{1}{3}$, která je zjevně splněna.

\[
\text{Řešení: } K = \{33\}.
\]
\end{examplebox}

To, co jsme doposud dělali, je velmi užitečné — umíme řešit rovnice, v nichž se
logaritmus vyskytuje samostatně, případně je jen vynásobený či upravený. Může však
nastat situace, kdy se v jediné rovnici objeví rovnou \emph{dva různé logaritmy},
navíc každý o jiném základu. V takových úlohách bývá neznámá ukryta hluboko
uvnitř vnořených logaritmů, takže je třeba postupovat pečlivě a krok za krokem.


Podívejme se na konkrétní příklad:

\[
\log_{0{,}5}\left(\log_4(x)\right) = -4.
\]

Na první pohled vypadá rovnici komplikovaně, ale ukáže se, že stačí správně rozbalit
logaritmy a použít odlogaritmování ve správném pořadí.

\begin{examplebox}
\textbf{Řešte rovnici:}
\[
\log_{0{,}5}(\log_4 x) = -4.
\]


\textbf{1. Určíme podmínky.}

Rovnice obsahuje vnořený logaritmus, proto musíme splnit:

1. vnitřní logaritmus $\log_4 x$ má smysl:
\[
x > 0,
\]

2. argument vnějšího logaritmu musí být kladný:
\[
\log_4 x > 0
\quad\Longleftrightarrow\quad
x > 1.
\]

Celková podmínka je tedy:
\[
x > 1.
\]

\textbf{2. Odlogaritmujeme vnější logaritmus.}

Použijeme ekvivalenci
\[
\log_{0{,}5}(A) = -4 \quad\Longleftrightarrow\quad A = (0{,}5)^{-4}.
\]

Proto:
\[
\log_4 x = (0{,}5)^{-4}.
\]

Vypočteme pravou stranu:
\[
(0{,}5)^{-4} = \left(\frac{1}{2}\right)^{-4} = 2^{4} = 16.
\]

Máme tedy:
\[
\log_4 x = 16.
\]

\textbf{3. Odlogaritmujeme i vnitřní logaritmus.}

Použijeme ekvivalenci $\log_4 x = 16 \Longleftrightarrow x = 4^{16}$:

\[
x = 4^{16}.
\]

\textbf{4. Ověříme podmínky.}

\[
4^{16} > 1 \quad\text{(platí)}.
\]

\[
\text{Řešení: } \{x\} = \{4^{16}\}.
\]

Poznámka: Zápis $4^{16}$ je naprosto validní a vhodný, neboť číselná hodnota je
extrémně velká ($4^{16} = 4\,294\,967\,296$).
\end{examplebox}



\subsection{Logaritmické nerovnice}

U logaritmických rovnic jsme převáděli logaritmy na exponenty a využívali tří
základních vzorců pro součin, podíl a mocninu uvnitř logaritmu. U logaritmických
\emph{nerovnic} je postup velmi podobný, ale přibývá jedna zásadní komplikace:
logaritmus může být \emph{rostoucí} nebo \emph{klesající} funkce podle svého
základu. A právě tato vlastnost rozhoduje o tom, zda se znaménko nerovnosti
při úpravách zachová, nebo obrátí.

\begin{definitionbox}
\textbf{Monotónnost logaritmické funkce.}

Nechť $a>0$, $a\neq 1$. Pak funkce $y = \log_a x$ má tyto vlastnosti:

\begin{enumerate}[label=\arabic*.]
  \item je \textbf{rostoucí}, pokud $a>1$,
  \item je \textbf{klesající}, pokud $0 < a < 1$.
\end{enumerate}
Z toho plyne, že:

\[
\log_a u < \log_a v \quad\Longleftrightarrow\quad
\begin{cases}
u < v, & a>1, \\
u > v, & 0<a<1.
\end{cases}
\]

\end{definitionbox}

\begin{notebox}
Připomínáme, že logaritmus je definován jen tehdy, když jeho argument je
kladný. U nerovnic proto bude prvním krokem vždy stanovení \textbf{podmínek
definice}.
\end{notebox}

\subsubsection*{Jednoduché nerovnice $\log_a(x) > c$}

Základní typ nerovnice řešíme úplným odlogaritmováním.

\begin{examplebox}
\textbf{Řešte nerovnici:}
\[
\log_2(x) > 3.
\]

\textbf{1. Podmínka:} musí platit $x>0$.

\textbf{2. Odlogaritmování:}
\[
\log_2(x) > 3 \quad\Longleftrightarrow\quad x > 2^3,
\]
protože $2>1$ je rostoucí logaritmus.

\[
x > 8.
\]

\textbf{Řešení: } $(8, +\infty)$.
\end{examplebox}

\begin{examplebox}
\textbf{Řešte nerovnici:}
\[
\log_{0{,}5}(x) > 3.
\]

\textbf{1. Podmínka:} $x > 0$.

\textbf{2. Odlogaritmování:} Základ je $0{,}5 < 1$, tedy logaritmus je
\emph{klesající} a znaménko se otočí:
\[
\log_{0{,}5}(x) > 3
\quad\Longleftrightarrow\quad
x < (0{,}5)^3.
\]

\[
x < \frac{1}{8}.
\]

\textbf{Řešení: } $(0,\, \tfrac18)$.
\end{examplebox}


\subsubsection*{Nerovnice ve tvaru $\log_a(\text{výraz}) > \log_a(\text{výraz})$}

Pokud mají oba logaritmy stejný základ, můžeme je „odstranit“ podle monotónnosti.

\begin{examplebox}
\textbf{Řešte nerovnici:}
\[
\log_3(2x+1) < \log_3(7-x).
\]

\textbf{1. Podmínky:} oba argumenty musí být kladné.
\[
2x+1 > 0 \quad\Longrightarrow\quad x > -\tfrac12,
\]
\[
7-x > 0 \quad\Longrightarrow\quad x < 7.
\]

\textbf{2. Protože $3>1$, logaritmus je rostoucí → znaménko se \textbf{nemění}:}
\[
2x+1 < 7 - x.
\]

\textbf{3. Vyřešíme lineární nerovnici:}
\[
3x < 6 \quad\Longrightarrow\quad x < 2.
\]

\textbf{4. Průnik s podmínkami:}
\[
x \in\left(-\tfrac12,\, 2\right).
\]

\textbf{Řešení: } $(-\tfrac12,\, 2)$.
\end{examplebox}


\subsubsection*{Nerovnice s více úpravami uvnitř logaritmu}

\begin{examplebox}
\textbf{Řešte nerovnici:}
\[
2\log_5(x-1) - \log_5(x+4) \ge 1.
\]

\textbf{1. Podmínky:}
\[
x-1 > 0 \Rightarrow x>1,\qquad
x+4>0 \Rightarrow x>-4.
\]
Celkově: $x>1$.

\textbf{2. Úpravy pomocí vzorců:}
\[
2\log_5(x-1) = \log_5((x-1)^2),
\]
\[
\log_5((x-1)^2) - \log_5(x+4)
= \log_5\!\left(\frac{(x-1)^2}{x+4}\right).
\]

Nerovnice přejde na:
\[
\log_5\!\left( \frac{(x-1)^2}{x+4} \right) \ge 1.
\]

\textbf{3. Odlogaritmování (logaritmus je rostoucí):}
\[
\frac{(x-1)^2}{x+4} \ge 5^1 = 5.
\]

\[
(x-1)^2 \ge 5(x+4).
\]

\textbf{4. Kvadratická nerovnice:}
\[
x^2 - 2x + 1 \ge 5x + 20,
\]
\[
x^2 - 7x -19 \ge 0.
\]

\[
\Delta = 49 + 76 = 125,
\quad
x = \frac{7 \pm \sqrt{125}}{2}.
\]

Znamení paraboly: otevřená nahoru → řešení je:
\[
x \le \frac{7 - \sqrt{125}}{2}
\quad\text{nebo}\quad
x \ge \frac{7 + \sqrt{125}}{2}.
\]

\textbf{5. Průnik s podmínkou $x>1$:}
Dolní interval odpadá, protože $\frac{7 - \sqrt{125}}{2} < 1$.

\[
\text{Řešení: } x \ge \frac{7 + \sqrt{125}}{2}.
\]
\end{examplebox}

\begin{notebox}
Nejčastější chyba: studenti zapomínají stanovit podmínky a zapíší mezi řešení i
hodnoty, pro které argument logaritmu není kladný. Logaritmické nerovnice se
proto vždy řeší ve dvou fázích:

\begin{enumerate}[label=\arabic*.]
\item Stanovení podmínek existence logaritmů.
\item Teprve poté řešení nerovnice.
\end{enumerate}

\end{notebox}



% -----------------------------------------------------------------
% ---------------------- P Ř Í K L A D Y --------------------------
% -----------------------------------------------------------------

\exercisepart
\setcounter{section}{1}

% =========================
% ÚLOHA 1 – Ověřte správnost tvrzení o logaritmech
% =========================
\begin{exercise}[Ověřte správnost následujících tvrzení]
    Pomocí definice logaritmu ověřte, zda jsou následující tvrzení pravdivá:
    \begin{tasks}(2)
        \task $\log_6 6 = 1$
        \task $\log_7 49 = 2$
        \task $\log 100 = 2$
        \task $\log_2 4 = 16$
        \task $\log_2 0{,}5 = -1$
        \task $\log_9 \sqrt{3} = 0{,}25$
        \task $\log_{27} \sqrt{27} = 0{,}125$
        \task $\log_5 25 = 3$
        \task $\log 0{,}01 = -2$
        \task $\log_3 81 = 4$
        \task $\log_4 16 = 3$
        \task $\log_{0{,}1} \frac{1}{10} = -1$
    \end{tasks}
\end{exercise}


% =========================
% ÚLOHA 2 – Vypočítejte hodnoty logaritmů
% =========================
\begin{exercise}[Vypočítejte hodnoty následujících logaritmů]
    Určete hodnotu každého z následujících logaritmů. Použijte definici dekadického logaritmu
    a základní vzorce pro práci s logaritmy.
    \begin{tasks}(2)
        \task $\log 10$
        \task $\log 100$
        \task $\log 10^{5}$
        \task $\log \sqrt{10}$
        \task $\log 10^{-3}$
        \task $\log 1$
        \task $\log 0{,}1$
        \task $\log 0{,}01$
    \end{tasks}
\end{exercise}


% =========================
% ÚLOHA 3 – Vypočítejte hodnoty logaritmických výrazů
% =========================
\begin{exercise}[Vypočítejte hodnoty následujících výrazů]
    Určete hodnotu výrazů využívajících základní ekvivalenci.
    \begin{tasks}(2)
        \task $5^{\log_5 2}$
        \task $10^{\log_{10} 1}$
        \task $8^{\,1+\log_{8} 5}$
        \task $3^{2\log_{3} 4}$
        \task $2^{\log_{2} 0{,}25}$
        \task $7^{\log_{7} 14 - \log_{7} 2}$
        \task $5^{\log_{5} 20 - \log_{5} 4}$
        \task $4^{\log_{4} 50 + \log_{4} 2}$
        \task $9^{\log_{9} 3^{3}}$
        \task $6^{\log_{6} 0{,}1}$
        \task $10^{\log 5 + \log 2}$
        \task $2^{\log_2 32 - \log_2 4}$
    \end{tasks}
\end{exercise}


% =========================
% ÚLOHA 4 – Jednoduché rovnice s logaritmem
% =========================
\begin{exercise}[Najděte všechny hodnoty proměnné]
    Najděte všechny hodnoty $x \in (0; +\infty)$, pro něž platí:
    \begin{tasks}(2)
        \task $\log_2 x = 0$
        \task $\log_2 x = 1$
        \task $\log_2 x = -3$
        \task $\log_2 x = -\frac{1}{5}$
        \task $\log_5 x = -3$
        \task $\log_5 x = 1$
        \task $\log_5 x = 0$
        \task $\log_5 x = -\frac{1}{3}$
        \task $\log_{10} x = 4$
        \task $\log_{10} x = -2$
        \task $\log_{10} x = 0$
        \task $\log_{10} x = 1$
    \end{tasks}
\end{exercise}


% =========================
% ÚLOHA 5 – Hledání vhodného základu logaritmu
% =========================
\begin{exercise}[Určete všechny hodnoty parametru $a$]
    Určete všechna reálná čísla $a$, pro něž platí:
    \begin{tasks}(2)
        \task $\log_a 4 = 2$
        \task $\log_a 0{,}0001 = -2$
        \task $\log_a 25 = 2$
        \task $\log_a 0{,}01 = -2$
        \task $\log_a 16 = 4$
        \task $\log_a 0{,}0001 = -4$
        \task $\log_a 9 = 2$
        \task $\log_a 0{,}09 = -2$
        \task $\log_a 81 = 4$
        \task $\log_a 0{,}0001 = -4$
    \end{tasks}
\end{exercise}


% =========================
% ÚLOHA 6 – Bonus: Náčrtky grafů netradičních funkcí
% =========================
\begin{exercise}[BONUS: Načrtněte grafy následujících funkcí]
    Načrtněte grafy těchto funkcí (uveďte také jejich definiční obory):
    \begin{tasks}(2)
        \task $y = \log_x 2$
        \task $y = 2^{\log_2 x}$
        \task $y = \log_x x$
        \task $y = x^{\log_x x}$
        \task $y = \log_x \left(\dfrac{1}{x}\right)$
        \task $y = \log_x \left(x^{\sqrt{x}}\right)$
        \task $y = \log_x \left(\log_x x\right)$
        \task $\star \; y = \log_x x \cdot \bigl(\log_x x^{\log_x (x^{\sqrt{x}})}\bigr)$
    \end{tasks}
\end{exercise}


% =========================
% ÚLOHA 7 – Upravte logaritmické výrazy
% =========================
\begin{exercise}[Upravte a vypočítejte hodnotu výrazů]
    Využijte vzorce pro součin, podíl a mocninu uvnitř logaritmu.
    \begin{tasks}(2)
        \task $4\log_6 3 \;+\; 5\log_6 2 \;-\; \log_6 12$
        \task $\log 1500 \;-\; \log 15$
        \task $\log_5 10 \;+\; \log_5 12{,}5$
        \task $-\log_3 4 \;+\; \log_3 6 \;-\; \dfrac{1}{4} \log_3 0{,}5$
        \task $\dfrac{2}{5} \left( 4\log_6 11 \;-\; \dfrac{1}{6} \log_6 4 \;+\; \log_6 5 \right)$
        \task $3\log_2 5 \;+\; 2\log_2 3 \;-\; \log_2 60$
        \task $\log_4 32 \;-\; 2\log_4 3 \;+\; \dfrac{1}{2} \log_4 9$
        \task $5\log_7 2 \;-\; 3\log_7 4 \;+\; \log_7 14$
        \task $2\log_3 5 \;+\; \log_3 9 \;-\; \log_3 20$
        \task $\log_2 40 \;-\; \log_2 5 \;+\; \dfrac{1}{2} \log_2 16$
        \task $-\log_8 27 \;+\; \log_8 9 \;+\; 2\log_8 6$
        \task $\dfrac{3}{2} \log_4 8 \;-\; \log_4 2 \;+\; \log_4 32$
    \end{tasks}
\end{exercise}


% =========================
% ÚLOHA 8 – Exponenciální pokles koncentrace DDT
% =========================
\begin{exercise}[Rozklad škodlivé látky DDT v potravním řetězci]

    \begin{minipage}[t]{0.68\textwidth}
        Látka DDT (dichlordifenyltrichlorethan), pro člověka velmi škodlivá, se hromadí 
        v potravním řetězci a dostává se tak do mléka a dalších potravin.
        Její koncentrace ve výši $5 \cdot 10^{-6}\,\%$ je v současné době ještě tolerovaná,
        do budoucna je však požadováno snížit ji na $2 \cdot 10^{-6}\,\%$.
        Používání DDT je dnes téměř ve všech státech zakázáno, ale chemický rozklad probíhá
        jen velmi pozvolna — poločas rozkladu je přibližně $30$ let.
    \end{minipage}
    \hfill
    \begin{minipage}[t]{0.28\textwidth}
        \begin{center}
          \smallskip
            \qrcode[height=2cm]{https://www.instagram.com/postrehy_z_wikipedie/p/C-aZd6siaIe/}\\[4pt]
            {\small Příspěvek k tématu na Instagramu, proč se zakázalo používat DDT.}
        \end{center}
    \end{minipage}

    \medskip

    \begin{tasks}(1)
        \task Předpokládejme, že pokles koncentrace lze popsat vztahem
        \[
            C(t) = C_0 \left(\frac{1}{2}\right)^{\frac{t}{30}},
        \]
        kde $C_0$ je počáteční koncentrace a $t$ čas v letech.  
        Zapište rovnici, kterou musí splňovat čas $t$, kdy koncentrace klesne z 
        $5 \cdot 10^{-6}\,\%$ na $2 \cdot 10^{-6}\,\%$.

        \task Vyřešte tuto rovnici a určete, za jak dlouho (na celé roky) 
        bude dosaženo požadované nižší koncentrace.
    \end{tasks}
\end{exercise}


% =========================
% ÚLOHA 9 – Složitější logaritmické rovnice
% =========================
\begin{exercise}[Řešte složitější logaritmické rovnice]
    Řešte následující rovnice, určete všechny reálné kořeny.
    \begin{tasks}(1)
        \task $\log_{2}(x+1) = 3$
        \task $4\log_{3}(2x-1) = 12$
        \task $\log_{5}(x-2) + \log_{5}(x+3) = 2$
        \task $2\log_{4}(x+1) - \log_{4}(x-3) = 2$
        \task $\log_{\frac12}(2-x) = -2$
        \task $\log_{4}(5x-4) = 2$
        \task $\log_{2}\log_{3}\log_{\frac12} x = 0$
        \task $\log_{\frac12}\log_{3}\bigl(1 + 20\log_{2} x\bigr) = -2$
        \task $\log_{2}\bigl[14 + 2\log_{7}\bigl(1 + 2\log_{\frac12} x\bigr)\bigr] = 4$
        \task $\log_{9}\Bigl\{3\log_{2}\bigl[1 + \log_{3}\bigl(1 - 2\log_{3} x\bigr)\bigr]\Bigr\} = 0{,}5$
        \task $\log_{3}\bigl[5 + \log_{2}(8 - x)\bigr] = 2$
        \task $\log_{\frac12}\Bigl\{4 + \log_{5}\bigl[10 + \log_{2}(x - 3)\bigr]\Bigr\} = -3$
        \task $\log_{5}\Bigl\{6 + \log_{2}\bigl[8 + \log_{3}(x - 1)\bigr]\Bigr\} = 2$
    \end{tasks}
\end{exercise}


% =========================
% ÚLOHA 10 – Řešte logaritmické rovnice
% =========================
\begin{exercise}[Řešte následující logaritmické rovnice]
    Určete všechna reálná čísla $x$, pro která mají rovnice smysl a která dané rovnice splňují.
    \begin{tasks}(2)
        \task $\log_3(x+5) = \log_3(2x-1)$
        \task $\log_5(x^2 - 17) = \log_5(x+3)$
        \task $\log_2(3x - 4) + \log_2(x + 1) = 3$
        \task $2\log_4(x - 2) - \log_4(x + 6) = 1$
        \task $\log_{0{,}5}(x + 1) + \log_{0{,}5}(x - 3) = 2$
        \task $3\log_2(x - 1) + \log_2(x + 3) = 5$
        \task $\log_7(x^2 - 4) - \log_7(x - 2) = 1$
        \task $2\log_3(x + 4) - \log_3(x - 1) = 2$
        \task $\log_{2}(x - 3) + \log_{2}(x + 5) = 4$
        \task $ \log_{5}(2x + 1) - \log_{5}(x - 1) = 1$
        \task $\log_{0{,}2}(x + 2) + \log_{0{,}2}(x - 4) = -3$
        \task $ \log_{4}(x^2 - 9) - \log_{4}(x + 3) = 1$
    \end{tasks}
\end{exercise}


% =========================
% ÚLOHA 11 – Další logaritmické rovnice
% =========================
\begin{exercise}[Řešte logaritmické rovnice s neznámou $x \in \mathbb{R}$]
    Najděte všechna reálná čísla $x$, pro která mají rovnice smysl a která je splňují.
    \begin{tasks}(2)
        \task $\log x = 2 - \log 5$
        \task $\log_{2}(x+14) + \log_{2}(x+2) = 6$
        \task $\log x = 2\log 5 + \log 4$
        \task $\log_{5}(x-1) + \log_{5}(x+3) = 2$
        \task $3\log_{2}(x+1) - \log_{2}(x-3) = 2$
        \task $\log_{3}(x+4) + \log_{9}(x-2) = 3$  
        \task $\log_{10}(x+3) - \log_{10} 2 = 1 - \log_{10}(x+2)$
        \task $\log_{4}(2x+5) - \log_{4}(x+1) = 1$
        \task $\log_{6}(x+1) + \log_{3} x = 1$
        \task $\log_{8}\sqrt{x+30} + \log_{8}\sqrt{x} = 1$
        \task \text{$\log_{10}\sqrt{x-4} + \log_{10}\sqrt{3x+1} = \log_{10} 40 - 1$}
        \task $\dfrac{2\log_{10}(x+2)}{\log_{10}(3x+10)} = 1$
        \task \text{$\log\sqrt{x} + \log\!\left(\dfrac{1}{x^{2}}\right) - \log x^{3} + \dfrac{11}{2}
               = \dfrac{\log x^{2}}{1 + \log 10}$}
    \end{tasks}
\end{exercise}


% =========================
% ÚLOHA 12 – Logaritmické rovnice s racionálními výrazy
% =========================
\begin{exercise}[Řešte logaritmické rovnice]
    Určete všechna reálná čísla $x$, pro která mají rovnice smysl a která dané rovnice splňují.

    \begin{tasks}(2)
        \task $\log_{2}\!\left(\frac{3 - x}{x + 3}\right) = -2$
        \task $\log_{5}\!\left(\frac{2x + 1}{x - 4}\right) = 1$
        \task $\log_{10}\!\left(\frac{x - 2}{x + 8}\right) = -1$
        \task $\log_{3}\!\left(\frac{6x - 2}{x - 3}\right) = 2$
        \task $\log_{\frac12}\!\left(\frac{x}{x + 14}\right) = \dfrac{\log 125}{\log 5}$
        \task $\log_{7}\!\left(\frac{x^{2} + 1}{x^{2} - 1}\right) = 1$
        \task $\log_{4}\!\left(\frac{x-2}{x+6}\right) + \log_{4}(x-1) = 1$
        \task $2\log_{3}(x+4) - \log_{3}(x^{2}-1) = \log_{3} 3$
        \task $\log_{0{,}5}\!\left(\frac{x+5}{2x-1}\right) - \log_{0{,}5}(x-3) = -2$
        \task $\log_{2}\!\left(\log_{3}(x+1)\right) = \log_{2}\!\left(\dfrac{1}{\,\log_{3}(x-2)}\right)$
        \task $\log_{5}\!\left(\frac{x^{2} - 4x + 4}{x^{2} - 9}\right) = 0$
        \task $ \log_{10}\!\left(\frac{2x + 6}{x - 3}\right) + \log_{10}(x - 5) = 1$
        \task $\log_{3}\!\left(\frac{x + 1}{x - 2}\right) + \log_{3}(x - 1) = 2$
        \task $\log_{2}\!\left(\frac{x^{2} - 5x + 6}{x^{2} - 4}\right) = 3$
    \end{tasks}
\end{exercise}


% =========================
% ÚLOHA 13 – Rovnice s logaritmy v zlomku
% =========================
\begin{exercise}[Řešte rovnice s neznámou]
    Řešte v množině reálných čísel $x$:
    \begin{tasks}(2)
        \task $\displaystyle \frac{\log_3(6x-2)}{\log_3(x-3)} = 2$
        \task $\displaystyle \frac{\log_5\left(x - \frac{1}{4}\right)}{\log_5\left(x + \frac{7}{2}\right)} = -1$
        \task $\displaystyle \frac{2\log(3x)}{\log(2-7x)} = 1$
        \task $\displaystyle \frac{\log x}{\log(x-2)} = \frac{\log 9}{\log 3}$
        \task $\displaystyle \frac{\log_2(x-1)}{\log_2(x-3)} = -1$
        \task $\displaystyle \frac{\log_3(2x+1)}{\log_3(x-2)} = 2$
        \task $\displaystyle \frac{2\log(3x-1)}{\log(x+4)} = 1$
        \task $\displaystyle \frac{\log(x^2-5x+6)}{\log(x-1)} = 2$
        \task $\displaystyle \frac{\log_2(x^2-1)}{\log_2(x-1)} = 3$
        \task $\displaystyle \frac{\log(x+1)}{\log(x-2)} = 2$
    \end{tasks}
\end{exercise}


% =========================
% ÚLOHA 14 – Řešte složitější logaritmické rovnice
% =========================
\begin{exercise}[Řešte složitější logaritmické rovnice]
    Řešte v množině reálných čísel $x$:
    \begin{tasks}(2)
        \task $\left(\log_{2} x\right)^{2} + 2\log_{2} x - 3 = 0$
        \task $4\log_{9} x \,(\log_{9} x - 1) = 2 + 3\log_{9} x$
        \task $\left(\log_{\frac12}(x+1)\right)^{2} + 5\log_{\frac12}(x+1) = 6$
        \task $\left(\log_{2}(x-1)\right)^{2} - 3\log_{2}(x-1) + 2 = 0$
        \task $2\left(\log_{3}(x+2)\right)^{2} - \log_{3}(x+2) - 1 = 0$
        \task $2\left(\log_{4}(x-3)\right)^{2} + 3\log_{4}(x-3) - 2 = 0$
        \task $\left(\log_{2}(x+5)\right)^{2} - 4\log_{2}(x+5) + 3 = 0$
        \task $3\left(\log_{7}(x-1)\right)^{2} - 2\log_{7}(x-1) - 1 = 0$
        \task $\left(\log_{10}(x+4)\right)^{2} + \log_{10}(x+4) - 6 = 0$        
        \task \text{$\left(\log_{5}(x^{2}-4)\right)^{2} - 3\log_{5}(x^{2}-4) + 2 = 0$}
        \task $\left(\log_{\frac13}(x-2)\right)^{2} + \log_{\frac13}(x-2) = \log_{\frac13} 9$
        \task $4\log_{3}(2x+1) + \log_{3}\sqrt{2x+1} = \frac{3}{2}\left(\log_{3}(2x+1)\right)^{2}$
    \end{tasks}
\end{exercise}


% =========================
% ÚLOHA 15 – Rovnice s logaritmem v exponentu
% =========================
\begin{exercise}[Řešte rovnice s logaritmem v exponentu]
    Řešte v množině reálných čísel $x$:
    \begin{tasks}(2)
        \task $x^{1+\log x} = 100$
        \task $\sqrt[5]{x^{\log_3 x}} = 243$
        \task $x^{\sqrt{x}} = (\sqrt{x})^{x}$
        \task $(\sqrt{x})^{\log_x (x-2)} = x$
        \task $x^{\log_2 x} = 16$
        \task $x^{\log_5 (x+4)} = 125$
        \task $(x-1)^{\log_x 9} = 27$
        \task $x^{\log_4 (x-3)} = 64$
    \end{tasks}
\end{exercise}


% =========================
% ÚLOHA 16 – Exponenciálně-logaritmické rovnice
% =========================
\begin{exercise}[Řešte exponenciálně-logaritmické rovnice]
    Řešte v množině reálných čísel $x$:
    \begin{tasks}(2)
        \task $x + \log_{2}\!\left(8 + 2^{x}\right) = 7$
        \task $x + \log_{2}\!\left(1 - 3\cdot 2^{x}\right) = x \log_{2} 4$
        \task $\log_{3}\!\left(4\cdot 2^{x}\right) + \log_{3} 3^{x+2} = \log_{3} 36^{x}$
        \task $\log_{5}\!\left(10\cdot 25^{x}\right) - \log_{5}\!\left(5^{x} + 25\right) = x + 1$
        \task $\log_{2}\!\left(5\cdot 2^{x} - 8\right) = x + 2$
        \task $\log_{3}\!\left(2\cdot 3^{x} - 3\right) = x$
        \task $\log_{4}\!\left(2^{x+1} + 8\right) = x$
        \task \text{$\log_{2} 3 + \log_{2} 4^{\,x+\sqrt{x}} = \log_{2}\!\left(2^{x+\sqrt{x}+1} + 4\right) + 2$}
        \task $\log_{5}\!\left(3\cdot 5^{x} - 10\right) + 1 = x$
        \task $\log_{2}\!\left(4 + 2^{x+1}\right) - x = 3$
    \end{tasks}
\end{exercise}


\setcounter{section}{2}
% =========================
% ÚLOHA 17 – Řešte logaritmické nerovnice
% =========================
\begin{exercise}[Řešte následující logaritmické nerovnice]
    Řešte v množině reálných čísel $x$:
    \begin{tasks}(2)
        \task $\log_{3}(x-1) > \log_{3} 3$
        \task $\log_{5}(2x+3) \leq 1$
        \task $\log_{10}(x-3) + \log_{10}(2x-1) < 0$
        \task $\log_{2}(x+1) \geq \log_{2}(5 - x)$
        \task $2\log_{4}(x-1) - \log_{4}(x+3) > 1$
        \task $\log_{0{,}5}(x-2) + \log_{0{,}5}(x-4) \leq -2$
        \task $\log_{3}(x+2) < \log_{3}(7 - x)$
        \task $3\log_{2}(x-1) + \log_{2}(x+3) \geq 5$
        \task $\log_{7}(x^{2} - 4) - \log_{7}(x - 2) < 1$
        \task $2\log_{3}(x + 4) - \log_{3}(x - 1) \leq 2$
        \task $\log_{2}(x - 3) + \log_{2}(x + 5) > 4$
        \task $\log_{5}(2x + 1) - \log_{5}(x - 1) \geq 1$
    \end{tasks}
\end{exercise}


% =========================
% ÚLOHA 18 – Logaritmické nerovnice
% =========================
\begin{exercise}[Řešte logaritmické nerovnice]
    Řešte v množině reálných čísel $x$:
    \begin{tasks}(2)
        \task $\log_{3}(2x - 1) \ge \log_{3}(4x + 3)$
        \task $\log_{\frac{1}{3}}(x + 4) \le \log_{\frac{1}{3}}(5x - 4)$
        \task $\log_{20}(x^{2} - 3x) \le \log_{20} x + \log_{20} 5$
        \task $\log_{0{,}2}(x - 3) > \log_{0{,}2} x - \log_{0{,}2} 2$
        \task $\log_{4}(x - 1) + \log_{4}(x + 3) > 2$
        \task $\log_{5}(x + 2) - \log_{5}(x - 3) < 1$
        \task $\log_{\frac{1}{2}}(x + 1) \ge \log_{\frac{1}{2}}(6 - x)$
        \task $\log_{7}(x - 4) + \log_{7}(x + 2) \le 2$
        \task $\log (2x - 5) - \log (x + 5) > -1$
        \task $2\log_{5}(x - 2) - \log_{5}(x + 1) \le 1$
        \task $\log_{\frac{1}{2}}(x - 1) + \log_{\frac{1}{2}}(x - 5) > -3$
        \task $\log_{3}(2x - 1) - \log_{3}(x + 2) \ge \log_{3} 3$
    \end{tasks}
\end{exercise}


\setcounter{section}{3}
% =========================
% ÚLOHA 19 – BONUS
% =========================
\begin{exercise}[BONUS: Dokažte vztah mezi logaritmy a exponenty]
    Je dáno, že $x = 2^{p}$ a $y = 4^{q}$, kde $p, q \in \mathbb{R}$.
    Dokažte, že platí
    \[
        \log_{2}\left(x^{3} \cdot y\right) = 3p + 2q.
    \]
    Svůj postup přehledně zdůvodněte pomocí vzorců pro logaritmus mocniny
    a vhodného zápisu čísla $4$ jako mocniny čísla $2$.
\end{exercise}


% =========================
% ÚLOHA 20 – BONUS
% =========================
\begin{exercise}[BONUS: Upravte rovnici exponenciální funkce]
    Exponenciální křivka má rovnici
    \[
        y = a b^{x}, \qquad x \in \mathbb{R},
    \]
    kde $a$ a $b$ jsou nenulové konstanty.

    Upravte tuto rovnici tak, aby neznámá $x$ byla vyjádřena jako subjekt
    (tj. vyjádřete ji explicitně v závislosti na ostatních veličinách).
    Konečný tvar výsledku napište pomocí logaritmů o základu $10$.
\end{exercise}


% =========================
% ÚLOHA 22 – BONUS
% =========================
\begin{exercise}[BONUS: Vyjádření neznámé bez logaritmů]
Je dáno, že reálné číslo $x$ splňuje logaritmickou rovnici
\[
    \log_{a} x = 2\bigl(\log_{a} k - \log_{a} 2\bigr),
\]
kde $k>0$, $a>0$ a $a \neq 1$.

\begin{enumerate}[label=\alph*)]
    \item Vyjádřete $x$ pomocí $k$ tak, aby výsledný tvar \emph{neobsahoval logaritmy}.
\end{enumerate}

Nyní předpokládejme, že číslo $x$ splňuje rovnici
\[
    \log_{x} (5y + 1) = 4 + \log_{x} 3,
\]
kde $x>0$, $x \neq 1$ a $y>0$, $y \neq 1$.

\begin{enumerate}[label=\alph*),start=2]
    \item Řešte tuto rovnici pro $y$ a vyjádřete $y$ pomocí $x$ tak, aby výsledný tvar \emph{neobsahoval logaritmy}.
\end{enumerate}
\end{exercise}


% =========================
% ÚLOHA 23 – Exponenciální rovnice (zaokrouhlení na 3 platné číslice)
% =========================
\begin{exercise}[BONUS: Řešte exponenciální rovnice]
    Řešte následující exponenciální rovnice. Výsledky zaokrouhlete na tři platné číslice.
    \begin{tasks}(2)
        \task $5^{2x - 1} = 4^{300}$
        \task $2^{y+1} = \dfrac{10}{2^{y}}$
    \end{tasks}
\end{exercise}


% =========================
% ÚLOHA 24 – BONUS: Úprava složitého logaritmického výrazu
% =========================
\begin{exercise}[BONUS: Zjednodušte logaritmický výraz]
    Zjednodušte co nejvíce následující logaritmický výraz a přitom přehledně ukažte všechny kroky úprav:
    \[
        \log\bigl(10 + 3\sqrt{10}\bigr)
        \;+\;
        \log\bigl(10 + \sqrt{90 + \sqrt{90}}\bigr)
        \;+\;
        \log\bigl(10 - \sqrt{90 + \sqrt{90}}\bigr).
    \]
    Ve svém postupu využijte vzorce pro logaritmus součinu a vhodné úpravy odmocnin.
\end{exercise}


% =========================
% ÚLOHA 25 – BONUS
% =========================
\begin{exercise}[BONUS: Najděte přesnou hodnotu $y$]
    Najděte přesnou hodnotu $y$, která splňuje rovnici
    \[
        2\log\!\left(\frac{x}{y}\right) - 1
        \;=\;
        \log\!\left(10x^{2} \cdot y\right),
        \qquad x \ne 0,\; y \ne 0.
    \]
    Výsledek vyjádřete přesně, bez použití desetinného zápisu.
\end{exercise}


% =========================
% ÚLOHA 26 – BONUS
% =========================
\begin{exercise}[BONUS: Určete parametry mocninné funkce]
    Body $(2,\,10)$ a $(6,\,100)$ leží na křivce dané rovnicí
    \[
        y = a x^{n},
        \qquad a \ne 0,\; n \ne 0.
    \]
    Určete hodnotu $a$ a hodnotu $n$ (zaokrouhlené na tři desetinná místa),
    které tuto křivku určují.
\end{exercise}


% =========================
% ÚLOHA 27 – BONUS
% =========================
\begin{exercise}[BONUS: Rovnice bez reálného řešení]
    Ukažte, že logaritmická rovnice
    \[
        \log_{x^{2} + 2}\!\bigl(2x^{4} - 2x^{3} + 7x^{2} - 2x + 5\bigr) = 2,
        \qquad x \in \mathbb{R},
    \]
    nemá žádné reálné řešení.

    Ve svém postupu nezapomeňte využít podmínky pro základ a argument
    logaritmu a vhodně upravit rovnici do tvaru mocninné rovnice.
\end{exercise}


\setcounter{section}{4}
% =========================
% ÚLOHA 28 – BONUS (\star)
% =========================
\begin{exercise}[$\star$ BONUS: Určete hodnotu logaritmu]
    Jsou dány rovnice
    \[
        \log(xy^{3}) = 1,
        \qquad
        \log(x^{2}y) = 1,
    \]
    kde předpokládáme $x>0$, $y>0$.
    
    Určete hodnotu výrazu
    \[
        \log(xy).
    \]

    Svůj postup přehledně zdůvodněte pomocí vzorců pro součin a mocninu uvnitř logaritmu.
\end{exercise}


% =========================
% ÚLOHA 29 – BONUS (\star)
% =========================
\begin{exercise}[$\star$ BONUS: Určete hodnotu součinu $abcd$]
    Nechť reálná čísla $a, b, c, d$ splňují následující vztahy:
    \[
        4^{a} = 5, \qquad
        5^{b} = 6, \qquad
        6^{c} = 7, \qquad
        7^{d} = 8.
    \]
    Určete hodnotu součinu
    \[
        a \cdot b \cdot c \cdot d.
    \]
    Svůj postup pečlivě zdůvodněte pomocí logaritmů a vlastností exponenciálních funkcí.
\end{exercise}


% =========================
% ÚLOHA 30 – BONUS
% =========================
\begin{exercise}[BONUS: Určete základ logaritmu]
    Řešení rovnice
    \[
        7^{\,x+7} = 8^{\,x}
    \]
    lze zapsat ve tvaru
    \[
        x = \log_{b}\!\left(7^{7}\right),
    \]
    kde $b$ je vhodně zvolený kladný základ logaritmu, $b \neq 1$.

    Určete hodnotu $b$.
\end{exercise}


% =========================
% ÚLOHA 31 – BONUS
% =========================
\begin{exercise}[BONUS: Vyjádření $x^{3}+y^{3}$ pomocí $10^{z}$]
    Nechť $S$ je množina všech uspořádaných trojic $(x,y,z)$ reálných čísel, pro které platí
    \[
        \log_{10}(x + y) = z
        \qquad\text{a}\qquad
        \log_{10}(x^{2} + y^{2}) = z + 1.
    \]

    Existují reálná čísla $a$ a $b$ taková, že pro všechny trojice $(x,y,z) \in S$ platí
    \[
        x^{3} + y^{3} = a \cdot 10^{3z} + b \cdot 10^{2z}.
    \]

    Určete hodnotu součtu
    \[
        a + b.
    \]
\end{exercise}


\end{document}